\section*{Overview}

The Chinese remainder theorem combines several modular constraints into one. In contests it turns ``find a number that
matches all these remainders'' into a constructive algorithm rather than a search.

\section*{When to Use It}

Use CRT when:

\begin{itemize}
  \item a value is constrained by several congruences,
  \item information is naturally split across different moduli,
  \item one large-modulus computation is easier after merging smaller ones.
\end{itemize}

\section*{Core Idea}

For coprime moduli, the system
\[
x \equiv r_1 \pmod{m_1}, \qquad x \equiv r_2 \pmod{m_2}
\]
has a unique solution modulo \(m_1 m_2\).

In the generalized version, two congruences are compatible iff
\[
r_1 \equiv r_2 \pmod{\gcd(m_1, m_2)}.
\]
If they are compatible, they can still be merged into one congruence modulo \(\mathrm{lcm}(m_1, m_2)\).

\section*{Key Insight}

CRT is really an alignment problem. Each congruence describes an arithmetic progression, and the theorem tells me when
those progressions overlap and how to describe the overlap compactly.

\section*{Operations / Main Technique}

\begin{itemize}
  \item combine two congruences at a time,
  \item use extended Euclid to solve the alignment step,
  \item iterate that merge across a whole system.
\end{itemize}

\paragraph{Worked example.}
The system
\[
x \equiv 2 \pmod 3, \qquad x \equiv 3 \pmod 5
\]
merges into
\[
x \equiv 8 \pmod{15}.
\]

\section*{Correctness Intuition}

When the residues agree modulo the gcd of the moduli, the two arithmetic progressions are synchronized often enough to
meet. Extended Euclid gives the coefficient needed to shift one progression onto the other. The merged modulus is the
period of that overlap.

\section*{Complexity Analysis}

Merging one pair of congruences takes \(O(\log m)\) because it relies on gcd and extended Euclid. Merging \(k\)
congruences is linear in \(k\) times that gcd cost.

\section*{Implementation}

The code supports the generalized pairwise merge, so it handles both coprime and merely compatible moduli. It returns a
boolean success flag together with the merged congruence.

\section*{Common Pitfalls}

\begin{itemize}
  \item Assuming the moduli are coprime when the statement does not say so.
  \item Forgetting to normalize the merged residue into \([0, m-1]\).
  \item Overflow when multiplying large moduli.
  \item Treating an inconsistent system as if it had a solution.
\end{itemize}

\section*{Variants / Extensions}

\begin{itemize}
  \item Garner's algorithm for reconstruction in mixed-radix form.
  \item Using CRT after several independent modular computations.
  \item Combining with modular inverse and extended Euclid in constructive number-theory problems.
\end{itemize}

\section*{Practice Problems}

\begin{itemize}
  \item Find the first time two or more periodic events align.
  \item Reconstruct a number from its residues.
  \item Merge congruence constraints in a simulation or scheduling problem.
\end{itemize}

\section*{References}

\begin{itemize}
  \item \href{https://cp-algorithms.com/algebra/chinese-remainder-theorem.html}{cp-algorithms: Chinese Remainder Theorem}.
  \item \href{https://codeforces.com/catalog}{Codeforces Catalog}.
\end{itemize}
